% blueprint_ladder_environments.tex -- the 38 theorem environments of planning/next_session_blueprint_node_set.md,
% chapters 2-9, drafted and rendered 2026-09-22 (every \lean{} name resolves; edges come from the walk, so no \uses here).
% Each block goes into blueprint/src/content.tex immediately BEFORE the anchor line named above it; one chapter per commit.
% Also in that pass: retitle the ch 5 and ch 7 sections "The theorem" -> "The theorems", and the \depgraphlegend clause
% for the shared box (both in the REPLACE items at the end). Not LaTeX input: a planning file.

% ===== chapter 2: insert before  \section{Example: MNIST MLP}
\begin{theorem}[Emitted train step is the certified step]
  \label{thm:mlp_fold}
  \lean{Proofs.MlpPoC.mlp_train_step_tied_certified}
  \leanok
  \Prove{} the six update operations \texttt{mlp\_train\_step.mlir} emits
  --- weights and biases of all three dense layers --- each denote
  \(\theta - \alpha\,\partial L/\partial\theta\): the output layer's against the
  gradient of the whole loss, the two hidden layers' against the certified
  per-layer VJP at the chain cotangent, with the loss cotangent
  \(\mathrm{softmax}(z) - \mathrm{onehot}(\ell)\) pinned to its proven graph.
\end{theorem}

\begin{proof}
\leanok
See \leandocref{Proofs.MlpPoC.mlp\_train\_step\_tied\_certified}.
\end{proof}

\begin{theorem}[SGD descends]
  \label{thm:mlp_sgd_descends}
  \lean{Proofs.mlp_input_sgd_descends, Proofs.mlp_input_float_sgd_descends}
  \leanok
  Under the pixel bound, an \(\eta\)-accurate gradient oracle, a margin at every
  ReLU of both hidden layers (the step stays inside the smooth cell), and the
  small-step and dominance conditions of Theorem~\ref{thm:linear_sgd_descends}:
  \Prove{} one SGD step on the input layer's weights \(W_0\) decreases the
  cross-entropy loss by at least \(\alpha\,\|\nabla L\|_2^2/2\).
  \texttt{mlp\_hidden\_sgd\_descends} is the same statement for \(W_1\); the
  binary32 twin supplies \(\eta\).
\end{theorem}

\begin{proof}
\leanok
See \leandocref{Proofs.mlp\_input\_sgd\_descends}.
\end{proof}

% ===== chapter 3: insert before  \section{Example: MNIST 2D CNN}
\begin{theorem}[Emitted convolution updates are the certified steps]
  \label{thm:cnn_fold}
  \lean{Proofs.CnnPoC.cnn_conv_tied_certified}
  \leanok
  \Prove{} the four convolution updates \texttt{cnn\_train\_step.mlir} emits
  --- \(W_2, b_2, W_1, b_1\) --- denote the certified SGD step at the real
  forward: each kernel update is \(W - \alpha\) times the conv weight-VJP at
  the chain cotangent, where the cotangent reaching \(\mathrm{conv}_2\) is the
  softmax cross-entropy gradient carried back through the dense head and the
  max pool, and \(\mathrm{conv}_1\)'s is that carried through \(\mathrm{conv}_2\).
  The dense head's updates tie as in Chapter~\ref{chap:mlp}.
\end{theorem}

\begin{proof}
\leanok
See \leandocref{Proofs.CnnPoC.cnn\_conv\_tied\_certified}.
\end{proof}

\begin{theorem}[SGD descends through the convolutions]
  \label{thm:cnn_sgd_descends}
  \lean{Proofs.cnn_conv2_sgd_descends, Proofs.cnn_conv2_float_sgd_descends}
  \leanok
  With a margin at the ReLUs after \(\mathrm{conv}_2\) and in the dense head, a
  strict maximum in every pooling window, an \(\eta\)-accurate gradient oracle
  and the small-step and dominance conditions: \Prove{} one SGD step on the
  second convolution's kernel \(W_2\) decreases the cross-entropy loss by at
  least \(\alpha\,\|\nabla L\|_2^2/2\). \texttt{cnn\_conv1\_sgd\_descends}
  reaches the first convolution; the binary32 twin supplies \(\eta\).
\end{theorem}

\begin{proof}
\leanok
See \leandocref{Proofs.cnn\_conv2\_sgd\_descends}.
\end{proof}

% ===== chapter 4: insert before  \section{Example: training dynamics on CIFAR}
\begin{theorem}[Emitted convolution updates are the certified steps]
  \label{thm:cifar_fold}
  \lean{Proofs.CifarPoC.cifar_conv_tied_certified}
  \leanok
  \Prove{} the eight convolution updates \texttt{cifar\_train\_step.mlir}
  emits --- kernels and biases of the two-stage CIFAR net's four convolutions
  --- denote the certified step at the real forward \texttt{cifarCnnForward},
  each at the chain cotangent carried back from the softmax cross-entropy
  gradient.
\end{theorem}

\begin{proof}
\leanok
See \leandocref{Proofs.CifarPoC.cifar\_conv\_tied\_certified}.
\end{proof}

\begin{theorem}[Emitted BatchNorm updates are the certified steps]
  \label{thm:cifar_bn_fold}
  \lean{Proofs.CifarBnPoC.bnSgdPairTied_holds}
  \leanok
  \Prove{} for a per-channel BatchNorm at any spatial shape and any
  cotangent \(c\), the two emitted updates denote the certified steps
  \(\gamma - \alpha \sum \pdiv(\mathrm{bn})\cdot c\) and
  \(\beta - \alpha \sum \pdiv(\mathrm{bn})\cdot c\) --- the scale and shift each
  take the gradient of Theorem~\ref{thm:bnHasVJP}'s forward.
\end{theorem}

\begin{proof}
\leanok
See \leandocref{Proofs.CifarBnPoC.bnSgdPairTied\_holds}.
\end{proof}

\begin{theorem}[The eight-convolution step is the certified step]
  \label{thm:cifar8_step_tie}
  \lean{Proofs.Cifar8PoC.cifar8_convs_tied_certified}
  \leanok
  \Prove{} all sixteen convolution updates of
  \texttt{cifar8\_train\_step.mlir} --- kernels and biases of \(W_1, \dots, W_8\)
  across the four stages --- denote the certified step at
  \texttt{cifarCnn8Forward}, cotangents chained back stage by stage from the
  loss gradient, and the dense head's updates are tied to the gradient of the
  whole loss.
\end{theorem}

\begin{proof}
\leanok
See \leandocref{Proofs.Cifar8PoC.cifar8\_convs\_tied\_certified}.
\end{proof}

\begin{theorem}[The eight-convolution BatchNorm step is the certified step]
  \label{thm:cifar8bn_step_tie}
  \lean{Proofs.Cifar8BnPoC.cifar8Bn_convbn_tied_certified}
  \leanok
  \Prove{} with BatchNorm after each of the eight convolutions, the
  twenty-four update pairs of \texttt{cifar8\_bn\_train\_step.mlir} --- kernel,
  bias, and BatchNorm \((\gamma, \beta)\) per layer --- denote the certified
  step at \texttt{cifarCnnBn8Forward}: the BatchNorm updates at the
  ReLU-masked cotangent, the convolution updates at the cotangent BatchNorm's
  input gradient passes down.
\end{theorem}

\begin{proof}
\leanok
See \leandocref{Proofs.Cifar8BnPoC.cifar8Bn\_convbn\_tied\_certified}.
\end{proof}

\begin{theorem}[SGD descends at the last convolution]
  \label{thm:cifar8_sgd_descends}
  \lean{Proofs.cifar8_lastConv_sgd_descends}
  \leanok
  With the first seven convolutions frozen, under the margin, small-step and
  dominance conditions: \Prove{} one SGD step on the last convolution's kernel
  \(W_8\) decreases the CIFAR-8 cross-entropy loss by at least
  \(\alpha\,\|\nabla L\|_2^2/2\) --- Theorem~\ref{thm:cnn_sgd_descends} at
  the factored forward \texttt{cifarCnn8Forward\_factor}.
\end{theorem}

\begin{proof}
\leanok
See \leandocref{Proofs.cifar8\_lastConv\_sgd\_descends}.
\end{proof}

% ===== chapter 5: insert before  \section{Example: ResNet-34 on Imagenette}
\begin{theorem}[ResNet-34 at batch \(N\) has the VJP]
  \label{thm:resnet34FullHasVJP}
  \lean{Proofs.resnet34ForwardBFullHasVJPAt_correct}
  \leanok
  \Assume{} the eighteen smooth-point bundles: no ReLU pre-activation of the
  stem, of any residual block, or of the pooled features sits at zero for the
  batch. \Prove{} the backward of the full-width ResNet-34 at training
  BatchNorm and batch \(N\) (\texttt{resnet34ForwardBFull}, \(3 \times 224
  \times 224\)) equals the \(\pdiv\)-contracted Jacobian of the forward at
  every input, cotangent and index. The statement is generic in \(N\) and in
  the weights.
\end{theorem}

\begin{proof}
\leanok
See \leandocref{Proofs.resnet34ForwardBFullHasVJPAt\_correct}.
\end{proof}

\begin{theorem}[ResNet-50 at batch \(N\) has the VJP]
  \label{thm:resnet50FullHasVJP}
  \lean{Proofs.resnet50ForwardBFullHasVJPAt_correct}
  \leanok
  The same for ResNet-50: with \(0 < q\) --- the image side is \(32q\), so
  \(q = 7\) is the 224-pixel net and \(q = 5\) the 160-pixel one --- and the
  eighteen bottleneck bundles, \Prove{} the backward of
  \texttt{resnet50ForwardBFull} at batch \(N\) equals its \(\pdiv\)
  Jacobian at every input, cotangent and index.
\end{theorem}

\begin{proof}
\leanok
See \leandocref{Proofs.resnet50ForwardBFullHasVJPAt\_correct}.
\end{proof}

\begin{theorem}[The rendered ResNet-50 backward is the VJP]
  \label{thm:resnet50_whole_back}
  \lean{Proofs.r50InputGradB_eq_r34B_full_vjp, Proofs.r50InputGradB_correct}
  \leanok
  \Prove{} the committed ResNet-50 input-gradient chain
  \texttt{r50InputGradB} --- the hand-written eighteen-stage reverse pass every
  \texttt{resnet50in\_*} artifact renders --- is the backward of the VJP above,
  and so equals
  \[
    \sum_j \pdiv(\texttt{resnet50ForwardBFull})\, x\, i\, j
  \cdot dy_j
  \]
  for every batch, cotangent and pixel.
\end{theorem}

\begin{proof}
\leanok
See \leandocref{Proofs.r50InputGradB\_eq\_r34B\_full\_vjp}.
\end{proof}

\begin{theorem}[Emitted ResNet-34 step is the certified step]
  \label{thm:resnet34_step_tie}
  \lean{Proofs.ResNet34TieB.r34_net_tiedB}
  \leanok
  \Prove{} all 110 emitted parameter gradients of the batched ResNet-34 step
  (\texttt{resnet34\_sgd\_train\_step.mlir} and \texttt{\_adam\_} at \(N = 32\),
  \texttt{resnet34in\_momdp64} at \(N = 64\)) are tied at the chain cotangent of
  the label-smoothed cross-entropy loss (\(\alpha_{\mathrm{ls}} = 0.1\), any
  target): each gradient node denotes the certified batch gradient
  \(\sum_n \mathrm{Jacobian} \cdot \mathrm{cotangent}\), generic in \(N\).
\end{theorem}

\begin{proof}
\leanok
See \leandocref{Proofs.ResNet34TieB.r34\_net\_tiedB}.
\end{proof}

\begin{theorem}[Emitted ResNet-50 step is the certified step]
  \label{thm:resnet50_step_tie}
  \lean{Proofs.ResNet50TieB.r50_net_tiedB, Proofs.ResNet50TieB.r50_lossCot_is_bce_grad}
  \leanok
  \Prove{} all 161 parameter gradients of the batched ResNet-50 step are tied
  at a loss cotangent \(g\) that is a binder, generic in \(N\) and \(q\);
  instantiated at the smoothed cross-entropy and at BCE-with-logits over
  \(N \cdot K\) (\texttt{r50\_lossCot\_is\_bce\_grad}) --- the loss of the
  76.66\% \texttt{resnet50in160\_lambaccdp8x64bce} run.
\end{theorem}

\begin{proof}
\leanok
See \leandocref{Proofs.ResNet50TieB.r50\_net\_tiedB}.
\end{proof}

\begin{theorem}[ResNet-34 seal: the backward is not degenerate]
  \label{thm:resnet34_seal}
  \lean{Proofs.R34FullBSeal.sealX_backward_nontrivial}
  \leanok
  \Prove{} at the structural weights \texttt{sealW} and input
  \texttt{sealX}, batch 2, there is a class \(j_0\) and a pixel \(i_0\) at
  which the backward of the full-width ResNet-34 is non-zero: the VJP theorem
  above is not satisfied by the zero map.
\end{theorem}

\begin{proof}
\leanok
See \leandocref{Proofs.R34FullBSeal.sealX\_backward\_nontrivial}.
\end{proof}

\begin{theorem}[ResNet-50 seal: the backward is not degenerate]
  \label{thm:resnet50_seal}
  \lean{Proofs.R50FullBSeal.sealX_backward_nontrivial}
  \leanok
  \Prove{} the same witness for ResNet-50, at both 224 and 160 pixels
  (\(0 < q \le 7\)).
\end{theorem}

\begin{proof}
\leanok
See \leandocref{Proofs.R50FullBSeal.sealX\_backward\_nontrivial}.
\end{proof}

\begin{theorem}[ResNet-34 data-parallel step is one device's step at \(R \cdot N\)]
  \label{thm:resnet34_sync_tie}
  \lean{Proofs.ResNet34SyncTieB.r34_net_syncTiedB}
  \leanok
  With \(R\) replicas of \(N\) images each and every BatchNorm synchronised:
  \Prove{} the mean over replicas of each of the 110 all-reduced gradient nodes
  of the data-parallel render (\texttt{resnet34in\_momdp64} at
  \(\texttt{replicas} > 1\), loss \(\div B\) per replica) equals the
  single-device node of Theorem~\ref{thm:resnet34_step_tie} at batch
  \(R \cdot N\) with loss \(\div (R \cdot B)\). The data-parallel step is one
  device's step on the global batch.
\end{theorem}

\begin{proof}
\leanok
See \leandocref{Proofs.ResNet34SyncTieB.r34\_net\_syncTiedB}.
\end{proof}

\begin{theorem}[ResNet-50 data-parallel step is one device's step at \(R \cdot N\)]
  \label{thm:resnet50_sync_tie}
  \lean{Proofs.ResNet50SyncTieB.r50_net_syncTiedB, Proofs.ResNet50SyncTieB.r50_net_syncTiedB_bce}
  \leanok
  \Prove{} the same for ResNet-50's 161 gradients, given that the shards are
  the scaled slices of one global batch; discharged at the smoothed
  cross-entropy and at BCE, with the committed divisors \(N \cdot K\) and
  \((R \cdot N) \cdot K\).
\end{theorem}

\begin{proof}
\leanok
See \leandocref{Proofs.ResNet50SyncTieB.r50\_net\_syncTiedB}.
\end{proof}

% ===== chapter 6: insert before  \section{Example: MobileNet V2 on Imagenette}
\begin{theorem}[MobileNetV2 at batch \(N\) has the VJP]
  \label{thm:mobilenetv2FullHasVJP}
  \lean{Proofs.mobilenetv2ForwardBFullHasVJPAt_correct}
  \leanok
  \Assume{} the nineteen ReLU6 bundles --- an expand clause and a depthwise
  clause per block, none after the residual add --- and the BatchNorm
  positivity conditions. \Prove{} the backward of the full-width MobileNetV2
  at training BatchNorm and batch \(N\) (\texttt{mobilenetv2ForwardBFull},
  seventeen blocks, \(3 \times 224 \times 224\)) equals the
  \(\pdiv\)-contracted Jacobian of the forward at every input, cotangent and
  index; generic in \(N\) and in the weights.
\end{theorem}

\begin{proof}
\leanok
See \leandocref{Proofs.mobilenetv2ForwardBFullHasVJPAt\_correct}.
\end{proof}

\begin{theorem}[MobileNetV4-Conv-M at batch \(N\) has the VJP]
  \label{thm:mobilenetv4FullHasVJP}
  \lean{Proofs.StableHLO.mobilenetv4ForwardBFullHasVJPAt_correct}
  \leanok
  The same for MobileNetV4-Conv-M (\texttt{mobilenetv4ForwardBFull}: the
  fused-IB stage and the universal inverted bottlenecks of
  \S~\ref{sec:mnv4_side_quest}) under one smooth-point bundle: \Prove{} its
  backward at batch \(N\) equals the \(\pdiv\) Jacobian at every input,
  cotangent and index.
\end{theorem}

\begin{proof}
\leanok
See \leandocref{Proofs.StableHLO.mobilenetv4ForwardBFullHasVJPAt\_correct}.
\end{proof}

\begin{theorem}[The rendered MobileNetV2 backward is the VJP]
  \label{thm:mobilenetv2_whole_back}
  \lean{Proofs.mnv2InputGradB_eq_mobilenetv2B_full_vjp, Proofs.mnv2InputGradB_correct}
  \leanok
  \Prove{} the committed MobileNetV2 input-gradient chain
  \texttt{mnv2InputGradB} is the backward of the twenty-one-stage VJP and
  equals
  \[
    \sum_j \pdiv(\texttt{mobilenetv2ForwardBFull})\, x\, i\, j \cdot
  dy_j
  \]
  for every batch, cotangent and pixel.
\end{theorem}

\begin{proof}
\leanok
See \leandocref{Proofs.mnv2InputGradB\_eq\_mobilenetv2B\_full\_vjp}.
\end{proof}

\begin{theorem}[The rendered MobileNetV4 backward is the VJP]
  \label{thm:mobilenetv4_whole_back}
  \lean{Proofs.mnv4InputGradB_eq_mnv4B_full_vjp, Proofs.mnv4InputGradB_correct}
  \leanok
  \Prove{} the same for \texttt{mnv4InputGradB} against the
  twenty-six-stage MobileNetV4 VJP, generic in \(N\) and in the number of
  classes.
\end{theorem}

\begin{proof}
\leanok
See \leandocref{Proofs.mnv4InputGradB\_eq\_mnv4B\_full\_vjp}.
\end{proof}

\begin{theorem}[Emitted MobileNetV2 step is the certified step]
  \label{thm:mobilenetv2_step_tie}
  \lean{Proofs.MobileNetV2TieB.mnv2_net_tiedB}
  \leanok
  \Prove{} all 158 emitted parameter gradients of the batched MobileNetV2
  step (\texttt{mobilenetv2\_adam\_train\_step.mlir} at \(N = 32\),
  \texttt{mobilenetv2in\_rmsdp64} at \(N = 64\)) are tied at the chain
  cotangent of the label-smoothed cross-entropy: each gradient node denotes
  the certified batch gradient, generic in \(N\).
\end{theorem}

\begin{proof}
\leanok
See \leandocref{Proofs.MobileNetV2TieB.mnv2\_net\_tiedB}.
\end{proof}

\begin{theorem}[Emitted MobileNetV4 step is the certified step]
  \label{thm:mobilenetv4_step_tie}
  \lean{Proofs.Mnv4TieB.mnv4_net_tiedB, Proofs.Mnv4TieB.mnv4_lossCot_is_smoothedCE_grad}
  \leanok
  \Prove{} all 233 parameter gradients of the batched MobileNetV4-Conv-M
  step (\texttt{mnv4in\_adamdp64}) are tied at a loss cotangent that is a
  binder, generic in \(N\) and in the number of classes;
  \texttt{mnv4\_lossCot\_is\_smoothedCE\_grad} instantiates it at the
  smoothed cross-entropy.
\end{theorem}

\begin{proof}
\leanok
See \leandocref{Proofs.Mnv4TieB.mnv4\_net\_tiedB}.
\end{proof}

\begin{theorem}[MobileNetV2 seal: the backward is not degenerate]
  \label{thm:mobilenetv2_seal}
  \lean{Proofs.Mnv2FullBSeal.sealX_backward_nontrivial}
  \leanok
  \Prove{} at structural weights and input, batch 2, there is a class and a
  pixel at which the backward of the full-width MobileNetV2 is non-zero ---
  thirty-five ReLU6 clauses, every one decided by the weights alone.
\end{theorem}

\begin{proof}
\leanok
See \leandocref{Proofs.Mnv2FullBSeal.sealX\_backward\_nontrivial}.
\end{proof}

\begin{theorem}[MobileNetV4 seal: the backward is not degenerate]
  \label{thm:mobilenetv4_seal}
  \lean{Proofs.Mnv4FullBSeal.sealX_backward_nontrivial}
  \leanok
  \Prove{} the same for MobileNetV4-Conv-M, fifty-four clauses; the fused
  swish makes the readout a swish gap over a grid-constant base rather than a
  product.
\end{theorem}

\begin{proof}
\leanok
See \leandocref{Proofs.Mnv4FullBSeal.sealX\_backward\_nontrivial}.
\end{proof}

\begin{theorem}[MobileNetV2 data-parallel step is one device's step at \(R \cdot N\)]
  \label{thm:mobilenetv2_sync_tie}
  \lean{Proofs.MobileNetV2SyncTieB.mnv2_net_syncTiedB}
  \leanok
  With \(R\) replicas of \(N\) images each and every BatchNorm synchronised:
  \Prove{} the mean over replicas of each of the 158 all-reduced gradient nodes
  of \texttt{mobilenetv2in\_rmsdp64} at \(\texttt{replicas} > 1\) (loss
  \(\div B\) per replica) equals the single-device node of
  Theorem~\ref{thm:mobilenetv2_step_tie} at batch \(R \cdot N\) with loss
  \(\div (R \cdot B)\).
\end{theorem}

\begin{proof}
\leanok
See \leandocref{Proofs.MobileNetV2SyncTieB.mnv2\_net\_syncTiedB}.
\end{proof}

\begin{theorem}[MobileNetV4 data-parallel step is one device's step at \(R \cdot N\)]
  \label{thm:mobilenetv4_sync_tie}
  \lean{Proofs.MobileNetV4SyncTieB.mnv4_net_syncTiedB, Proofs.MobileNetV4SyncTieB.mnv4_net_syncTiedB_smoothedCE}
  \leanok
  \Prove{} the same for MobileNetV4's 233 gradients, given that the shards
  are the scaled slices of one global batch; discharged at the smoothed
  cross-entropy.
\end{theorem}

\begin{proof}
\leanok
See \leandocref{Proofs.MobileNetV4SyncTieB.mnv4\_net\_syncTiedB}.
\end{proof}

% ===== chapter 7: insert before  \section{Example: EfficientNet-B0 on Imagenette}
\begin{theorem}[EfficientNet-B0 at batch \(N\): VJP and graph]
  \label{thm:efficientnetFullHasVJP}
  \lean{Proofs.efficientnetForwardBFullHasVJP_correct, Proofs.StableHLO.efficientnetFwdGraphBFull_faithful}
  \leanok
  \Assume{} the fifty BatchNorm positivities. \Prove{} the batched
  EfficientNet-B0 --- all sixteen MBConv blocks with squeeze-and-excitation
  and swish, at training BatchNorm and batch \(N\) --- has the VJP: its
  backward equals the \(\pdiv\)-contracted Jacobian of
  \texttt{efficientnetForwardBFull} at every input, cotangent and pixel;
  and the emitted forward graph denotes exactly that forward
  (\texttt{efficientnetFwdGraphBFull\_faithful}).
\end{theorem}

\begin{proof}
\leanok
See \leandocref{Proofs.efficientnetForwardBFullHasVJP\_correct}.
\end{proof}

\begin{theorem}[The rendered backward is the VJP]
  \label{thm:efficientnet_whole_back}
  \lean{Proofs.efficientnetInputGradBFull_correct}
  \leanok
  \Prove{} the hand-written sixteen-block reverse chain
  \texttt{efficientnetInputGradBFull} equals
  \[
    \sum_j \pdiv(\texttt{efficientnetForwardBFull})\, x\, i\, j \cdot dy_j
  \]
  for every batch, cotangent and pixel.
\end{theorem}

\begin{proof}
\leanok
See \leandocref{Proofs.efficientnetInputGradBFull\_correct}.
\end{proof}

\begin{theorem}[Emitted train step is the certified step]
  \label{thm:efficientnet_step_tie}
  \lean{Proofs.EnetTiePoCG.efficientnet_net_tiedG}
  \leanok
  \Prove{} the 262 parameter gradients of the batched EfficientNet-B0 step
  (\texttt{efficientnet\_adam\_train\_step.mlir} and every
  \texttt{efficientnetin\_*} artifact) are tied at the chain cotangent of the
  label-smoothed cross-entropy, the target a binder: each raw gradient node
  denotes the certified batch gradient through swish, the SE gate's fan-in
  and the batch-BatchNorm backward. The AdamW, RMSProp, EMA and
  data-parallel tails all consume that node.
\end{theorem}

\begin{proof}
\leanok
See \leandocref{Proofs.EnetTiePoCG.efficientnet\_net\_tiedG}.
\end{proof}

\begin{theorem}[The data-parallel step is one device's step at \(R \cdot N\)]
  \label{thm:efficientnet_sync_tie}
  \lean{Proofs.EnetSyncTieG.efficientnet_net_syncTiedG}
  \leanok
  With \(R\) replicas and all 49 BatchNorms synchronised: \Prove{} the mean
  over replicas of each of the 213 all-reduced gradient nodes (loss \(\div B\)
  per replica) equals the single-device node of
  Theorem~\ref{thm:efficientnet_step_tie} at batch \(R \cdot N\), loss
  \(\div (R \cdot B)\).
\end{theorem}

\begin{proof}
\leanok
See \leandocref{Proofs.EnetSyncTieG.efficientnet\_net\_syncTiedG}.
\end{proof}

% ===== chapter 8: insert before  \section{Example: ConvNeXt-T on Imagenette}
\begin{theorem}[The rendered backward is the VJP]
  \label{thm:convnext_whole_back}
  \lean{Proofs.convnextInputGradB_correct, Proofs.convnextImagenetInputGradB_eq_vjp}
  \leanok
  \Assume{} the twenty-three LayerNorm positivities. \Prove{} the
  twelve-stage ConvNeXt-T input-gradient chain \texttt{convnextInputGradB}
  equals the backward of the batched forward,
  \[
    \sum_j \pdiv(\mathrm{batchMap}\, B\; \texttt{convNextForwardTCh})\, x\, i\, j
  \cdot dy_j
  \]
  , for every batch, every number of classes and every cotangent;
  \texttt{convnextImagenetInputGradB\_eq\_vjp} is the 1000-class instance.
\end{theorem}

\begin{proof}
\leanok
See \leandocref{Proofs.convnextInputGradB\_correct}.
\end{proof}

\begin{theorem}[Emitted train step is the certified step]
  \label{thm:convnext_step_tie}
  \lean{Proofs.CnxTiePoCGB.cnx_net_tiedGB}
  \leanok
  \Prove{} the 182 parameter gradients of the batched ConvNeXt-T step
  (\([3, 3, 9, 3]\) blocks, channel-LayerNorm at twenty-two sites, per-channel
  layer scale, GELU) are tied at the chain cotangent of the smoothed
  cross-entropy loss, with the batch and the number of classes binders: each
  raw gradient node of \texttt{convnext\_adam\_train\_step.mlir} and of every
  \texttt{convnextin\_*} artifact denotes the certified batch gradient.
\end{theorem}

\begin{proof}
\leanok
See \leandocref{Proofs.CnxTiePoCGB.cnx\_net\_tiedGB}.
\end{proof}

% ===== chapter 9: insert before  \begin{theorem}[ViT-Tiny, the whole network]
\begin{theorem}[The depth-\(k\) tower has a VJP]
  \label{thm:vitForwardKVHasVJP}
  \lean{Proofs.vitForwardKVHasVJP_correct}
  \leanok
  \Assume{} \(\varepsilon > 0\). \Prove{} the depth-\(k\) transformer tower
  with distinct parameters per block --- patch embedding, CLS and position
  tables, \(k\) blocks of multi-head self-attention and MLP sublayers under
  vector LayerNorm, the classifier head --- has a global VJP, and its backward
  equals the \(\pdiv\)-contracted Jacobian of \texttt{vitForwardKV} at every
  input, cotangent and index, for every \(k\).
\end{theorem}

\begin{proof}
\leanok
See \leandocref{Proofs.vitForwardKVHasVJP\_correct}.
\end{proof}

\begin{theorem}[The depth-\(k\) graph denotes the tower]
  \label{thm:vitFwdGraphKMHV_faithful}
  \lean{Proofs.StableHLO.vitFwdGraphKMHV_faithful}
  \leanok
  \Prove{} for every depth \(k\) the emitted forward graph
  \texttt{vitFwdGraphKMHV} denotes \texttt{vitForwardKV}: the rendered ViT is
  the tower above.
\end{theorem}

\begin{proof}
\leanok
See \leandocref{Proofs.StableHLO.vitFwdGraphKMHV\_faithful}.
\end{proof}

% ===== chapter 9: insert before  \section{Example: ViT-Tiny on Imagenette}
\begin{theorem}[The rendered backward is the VJP]
  \label{thm:vit_whole_back}
  \lean{Proofs.vitTinyInputGradB_eq_vitTiny_vjp, Proofs.vitInputGradKB_correct}
  \leanok
  \Prove{} at ViT-Tiny's dimensions the batched input-gradient chain
  \texttt{vitInputGradKB} equals the backward of the batched tower
  \(\mathrm{batchMap}\, B\; \texttt{vitForwardKV}\) for every \(B\), cotangent
  and index, with only \(\varepsilon > 0\) assumed; the general-depth reading
  \texttt{vitInputGradKB\_correct} leaves the number of classes free.
\end{theorem}

\begin{proof}
\leanok
See \leandocref{Proofs.vitTinyInputGradB\_eq\_vitTiny\_vjp}.
\end{proof}

\begin{theorem}[Emitted train step is the certified step]
  \label{thm:vit_step_tie}
  \lean{Proofs.ViTTiePoCGB.vit_net_tiedGB}
  \leanok
  \Prove{} the 200 parameter gradients of the batched ViT-Tiny step (three
  heads of 64, depth 12, \(16 \times 16\) patches) are tied at the chain
  cotangent of the smoothed cross-entropy loss, with the batch and the number
  of classes binders; the CLS gradient is summed over the batch inside the
  node's denotation. Every \texttt{vitin\_*} artifact renders from this chain.
\end{theorem}

\begin{proof}
\leanok
See \leandocref{Proofs.ViTTiePoCGB.vit\_net\_tiedGB}.
\end{proof}

% ===== REPLACE items (old -> new), applied once each
% --- old:
%   \section{The theorem}
%   
%   \begin{center}
\begin{tikzpicture}[x=1pt,y=1pt, every node/.style={inner sep=1.50pt, font=\fontsize{7.00}{8.40}\selectfont\ttfamily}]
  \node[draw=green!50!black, fill=green!15, rectangle, rounded corners=4.0pt, align=center, minimum width=73.2pt, minimum height=12.7pt] at (126.7pt,8.9pt) {\hyperref[thm:residual_has_vjp]{residual\_has\_vjp}};
  \node[draw=green!50!black, fill=green!15, rectangle, rounded corners=4.0pt, align=center, minimum width=56.4pt, minimum height=21.5pt] at (126.7pt,213.9pt) {\hyperref[thm:resnet34_full_has_vjp]{resnet34\_}\\\hyperref[thm:resnet34_full_has_vjp]{full\_has\_vjp}};
  \node[draw=green!50!black, fill=green!15, rectangle, rounded corners=4.0pt, align=center, minimum width=56.4pt, minimum height=21.5pt] at (126.7pt,180.9pt) {\hyperref[thm:resnet50_full_has_vjp]{resnet50\_}\\\hyperref[thm:resnet50_full_has_vjp]{full\_has\_vjp}};
  \node[draw=green!50!black, fill=green!15, rectangle, rounded corners=4.0pt, align=center, minimum width=85.8pt, minimum height=12.7pt] at (126.7pt,152.9pt) {\hyperref[thm:resnet50_whole_back]{resnet50\_whole\_back}};
  \node[draw=green!50!black, fill=green!15, rectangle, rounded corners=4.0pt, align=center, minimum width=77.4pt, minimum height=12.7pt] at (126.7pt,128.9pt) {\hyperref[thm:resnet34_step_tie]{resnet34\_step\_tie}};
  \node[draw=green!50!black, fill=green!15, rectangle, rounded corners=4.0pt, align=center, minimum width=77.4pt, minimum height=12.7pt] at (126.7pt,241.9pt) {\hyperref[thm:resnet50_step_tie]{resnet50\_step\_tie}};
  \node[draw=green!50!black, fill=green!15, rectangle, rounded corners=4.0pt, align=center, minimum width=60.6pt, minimum height=12.7pt] at (126.7pt,104.9pt) {\hyperref[thm:resnet34_seal]{resnet34\_seal}};
  \node[draw=green!50!black, fill=green!15, rectangle, rounded corners=4.0pt, align=center, minimum width=60.6pt, minimum height=12.7pt] at (126.7pt,80.9pt) {\hyperref[thm:resnet50_seal]{resnet50\_seal}};
  \node[draw=green!50!black, fill=green!15, rectangle, rounded corners=4.0pt, align=center, minimum width=77.4pt, minimum height=12.7pt] at (126.7pt,56.9pt) {\hyperref[thm:resnet34_sync_tie]{resnet34\_sync\_tie}};
  \node[draw=green!50!black, fill=green!15, rectangle, rounded corners=4.0pt, align=center, minimum width=77.4pt, minimum height=12.7pt] at (126.7pt,32.9pt) {\hyperref[thm:resnet50_sync_tie]{resnet50\_sync\_tie}};
  \node[draw=gray!70, dashed, rounded corners=2pt, minimum width=52.2pt, minimum height=17.8pt, text=gray!60!black, align=center, font=\fontsize{5.50}{6.60}\selectfont\ttfamily] at (29.4pt,8.9pt) {\hyperref[thm:biPath_has_vjp]{biPath\_has\_vjp}\\(ch 1)};
  \node[draw=gray!70, dashed, rounded corners=2pt, minimum width=55.5pt, minimum height=17.8pt, text=gray!60!black, align=center, font=\fontsize{5.50}{6.60}\selectfont\ttfamily] at (29.4pt,241.9pt) {\hyperref[thm:pdiv_finset_sum]{pdiv\_finset\_sum}\\(ch 1)};
  \node[draw=gray!70, dashed, rounded corners=2pt, minimum width=58.8pt, minimum height=17.8pt, text=gray!60!black, align=center, font=\fontsize{5.50}{6.60}\selectfont\ttfamily] at (29.4pt,116.9pt) {20 statements of\\ch 1, 2, 3, 4, 9};
  \draw[->, >=stealth, dashed, gray!60, line width=0.40pt] (55.9pt,8.9pt) .. controls (62.9pt,8.9pt) and (70.7pt,8.9pt) .. (78.4pt,8.9pt) -- (89.7pt,8.9pt);
  \draw[->, >=stealth, dashed, gray!60, line width=0.40pt] (57.5pt,241.9pt) .. controls (63.3pt,241.9pt) and (69.7pt,241.9pt) .. (76.1pt,241.9pt) -- (87.6pt,241.9pt);
  \draw[->, >=stealth, dashed, gray!60, line width=0.40pt] (34.3pt,126.1pt) .. controls (41.7pt,142.5pt) and (59.2pt,176.8pt) .. (83.8pt,196.9pt) .. controls (85.3pt,198.1pt) and (86.8pt,199.2pt) .. (88.4pt,200.2pt) -- (98.2pt,205.5pt);
  \draw[->, >=stealth, dashed, gray!60, line width=0.40pt] (38.7pt,126.1pt) .. controls (48.5pt,136.7pt) and (65.9pt,154.0pt) .. (83.8pt,164.9pt) .. controls (85.2pt,165.7pt) and (86.7pt,166.6pt) .. (88.2pt,167.4pt) -- (98.4pt,172.1pt);
  \draw[->, >=stealth, dashed, gray!60, line width=0.40pt] (49.4pt,126.2pt) .. controls (59.5pt,131.0pt) and (72.2pt,136.6pt) .. (83.8pt,140.9pt) .. controls (85.6pt,141.5pt) and (87.4pt,142.2pt) .. (89.2pt,142.8pt) -- (99.8pt,146.1pt);
  \draw[->, >=stealth, dashed, gray!60, line width=0.40pt] (59.3pt,120.5pt) .. controls (64.7pt,121.2pt) and (70.6pt,121.9pt) .. (76.5pt,122.7pt) -- (87.8pt,124.1pt);
  \draw[->, >=stealth, dashed, gray!60, line width=0.40pt] (59.3pt,113.2pt) .. controls (67.2pt,112.2pt) and (76.1pt,111.1pt) .. (84.6pt,110.1pt) -- (96.0pt,108.6pt);
  \draw[->, >=stealth, dashed, gray!60, line width=0.40pt] (49.4pt,107.5pt) .. controls (59.5pt,102.8pt) and (72.2pt,97.1pt) .. (83.8pt,92.9pt) .. controls (85.6pt,92.2pt) and (87.4pt,91.6pt) .. (89.2pt,91.0pt) -- (99.8pt,87.7pt);
  \draw[->, >=stealth, dashed, gray!60, line width=0.40pt] (38.2pt,107.7pt) .. controls (47.8pt,96.9pt) and (65.3pt,79.1pt) .. (83.8pt,68.9pt) .. controls (84.3pt,68.6pt) and (84.7pt,68.4pt) .. (85.2pt,68.1pt) -- (95.7pt,63.6pt);
  \draw[->, >=stealth, dashed, gray!60, line width=0.40pt] (34.6pt,107.9pt) .. controls (41.5pt,94.0pt) and (56.6pt,67.2pt) .. (84.5pt,46.0pt) -- (93.6pt,39.7pt);
\end{tikzpicture}
\end{center}

% --- new:
%   \section{The theorems}
%   
%   \begin{center}
\begin{tikzpicture}[x=1pt,y=1pt, every node/.style={inner sep=1.50pt, font=\fontsize{7.00}{8.40}\selectfont\ttfamily}]
  \node[draw=green!50!black, fill=green!15, rectangle, rounded corners=4.0pt, align=center, minimum width=73.2pt, minimum height=12.7pt] at (126.7pt,8.9pt) {\hyperref[thm:residual_has_vjp]{residual\_has\_vjp}};
  \node[draw=green!50!black, fill=green!15, rectangle, rounded corners=4.0pt, align=center, minimum width=56.4pt, minimum height=21.5pt] at (126.7pt,213.9pt) {\hyperref[thm:resnet34_full_has_vjp]{resnet34\_}\\\hyperref[thm:resnet34_full_has_vjp]{full\_has\_vjp}};
  \node[draw=green!50!black, fill=green!15, rectangle, rounded corners=4.0pt, align=center, minimum width=56.4pt, minimum height=21.5pt] at (126.7pt,180.9pt) {\hyperref[thm:resnet50_full_has_vjp]{resnet50\_}\\\hyperref[thm:resnet50_full_has_vjp]{full\_has\_vjp}};
  \node[draw=green!50!black, fill=green!15, rectangle, rounded corners=4.0pt, align=center, minimum width=85.8pt, minimum height=12.7pt] at (126.7pt,152.9pt) {\hyperref[thm:resnet50_whole_back]{resnet50\_whole\_back}};
  \node[draw=green!50!black, fill=green!15, rectangle, rounded corners=4.0pt, align=center, minimum width=77.4pt, minimum height=12.7pt] at (126.7pt,128.9pt) {\hyperref[thm:resnet34_step_tie]{resnet34\_step\_tie}};
  \node[draw=green!50!black, fill=green!15, rectangle, rounded corners=4.0pt, align=center, minimum width=77.4pt, minimum height=12.7pt] at (126.7pt,241.9pt) {\hyperref[thm:resnet50_step_tie]{resnet50\_step\_tie}};
  \node[draw=green!50!black, fill=green!15, rectangle, rounded corners=4.0pt, align=center, minimum width=60.6pt, minimum height=12.7pt] at (126.7pt,104.9pt) {\hyperref[thm:resnet34_seal]{resnet34\_seal}};
  \node[draw=green!50!black, fill=green!15, rectangle, rounded corners=4.0pt, align=center, minimum width=60.6pt, minimum height=12.7pt] at (126.7pt,80.9pt) {\hyperref[thm:resnet50_seal]{resnet50\_seal}};
  \node[draw=green!50!black, fill=green!15, rectangle, rounded corners=4.0pt, align=center, minimum width=77.4pt, minimum height=12.7pt] at (126.7pt,56.9pt) {\hyperref[thm:resnet34_sync_tie]{resnet34\_sync\_tie}};
  \node[draw=green!50!black, fill=green!15, rectangle, rounded corners=4.0pt, align=center, minimum width=77.4pt, minimum height=12.7pt] at (126.7pt,32.9pt) {\hyperref[thm:resnet50_sync_tie]{resnet50\_sync\_tie}};
  \node[draw=gray!70, dashed, rounded corners=2pt, minimum width=52.2pt, minimum height=17.8pt, text=gray!60!black, align=center, font=\fontsize{5.50}{6.60}\selectfont\ttfamily] at (29.4pt,8.9pt) {\hyperref[thm:biPath_has_vjp]{biPath\_has\_vjp}\\(ch 1)};
  \node[draw=gray!70, dashed, rounded corners=2pt, minimum width=55.5pt, minimum height=17.8pt, text=gray!60!black, align=center, font=\fontsize{5.50}{6.60}\selectfont\ttfamily] at (29.4pt,241.9pt) {\hyperref[thm:pdiv_finset_sum]{pdiv\_finset\_sum}\\(ch 1)};
  \node[draw=gray!70, dashed, rounded corners=2pt, minimum width=58.8pt, minimum height=17.8pt, text=gray!60!black, align=center, font=\fontsize{5.50}{6.60}\selectfont\ttfamily] at (29.4pt,116.9pt) {20 statements of\\ch 1, 2, 3, 4, 9};
  \draw[->, >=stealth, dashed, gray!60, line width=0.40pt] (55.9pt,8.9pt) .. controls (62.9pt,8.9pt) and (70.7pt,8.9pt) .. (78.4pt,8.9pt) -- (89.7pt,8.9pt);
  \draw[->, >=stealth, dashed, gray!60, line width=0.40pt] (57.5pt,241.9pt) .. controls (63.3pt,241.9pt) and (69.7pt,241.9pt) .. (76.1pt,241.9pt) -- (87.6pt,241.9pt);
  \draw[->, >=stealth, dashed, gray!60, line width=0.40pt] (34.3pt,126.1pt) .. controls (41.7pt,142.5pt) and (59.2pt,176.8pt) .. (83.8pt,196.9pt) .. controls (85.3pt,198.1pt) and (86.8pt,199.2pt) .. (88.4pt,200.2pt) -- (98.2pt,205.5pt);
  \draw[->, >=stealth, dashed, gray!60, line width=0.40pt] (38.7pt,126.1pt) .. controls (48.5pt,136.7pt) and (65.9pt,154.0pt) .. (83.8pt,164.9pt) .. controls (85.2pt,165.7pt) and (86.7pt,166.6pt) .. (88.2pt,167.4pt) -- (98.4pt,172.1pt);
  \draw[->, >=stealth, dashed, gray!60, line width=0.40pt] (49.4pt,126.2pt) .. controls (59.5pt,131.0pt) and (72.2pt,136.6pt) .. (83.8pt,140.9pt) .. controls (85.6pt,141.5pt) and (87.4pt,142.2pt) .. (89.2pt,142.8pt) -- (99.8pt,146.1pt);
  \draw[->, >=stealth, dashed, gray!60, line width=0.40pt] (59.3pt,120.5pt) .. controls (64.7pt,121.2pt) and (70.6pt,121.9pt) .. (76.5pt,122.7pt) -- (87.8pt,124.1pt);
  \draw[->, >=stealth, dashed, gray!60, line width=0.40pt] (59.3pt,113.2pt) .. controls (67.2pt,112.2pt) and (76.1pt,111.1pt) .. (84.6pt,110.1pt) -- (96.0pt,108.6pt);
  \draw[->, >=stealth, dashed, gray!60, line width=0.40pt] (49.4pt,107.5pt) .. controls (59.5pt,102.8pt) and (72.2pt,97.1pt) .. (83.8pt,92.9pt) .. controls (85.6pt,92.2pt) and (87.4pt,91.6pt) .. (89.2pt,91.0pt) -- (99.8pt,87.7pt);
  \draw[->, >=stealth, dashed, gray!60, line width=0.40pt] (38.2pt,107.7pt) .. controls (47.8pt,96.9pt) and (65.3pt,79.1pt) .. (83.8pt,68.9pt) .. controls (84.3pt,68.6pt) and (84.7pt,68.4pt) .. (85.2pt,68.1pt) -- (95.7pt,63.6pt);
  \draw[->, >=stealth, dashed, gray!60, line width=0.40pt] (34.6pt,107.9pt) .. controls (41.5pt,94.0pt) and (56.6pt,67.2pt) .. (84.5pt,46.0pt) -- (93.6pt,39.7pt);
\end{tikzpicture}
\end{center}


% --- old:
%   \section{The theorem}
%   
%   \begin{center}
\begin{tikzpicture}[x=1pt,y=1pt, every node/.style={inner sep=1.50pt, font=\fontsize{7.00}{8.40}\selectfont\ttfamily}]
  \node[draw=green!50!black, fill=green!15, rectangle, rounded corners=4.0pt, align=center, minimum width=69.0pt, minimum height=12.7pt] at (80.9pt,6.4pt) {\hyperref[thm:seBlock_has_vjp]{seBlock\_has\_vjp}};
  \node[draw=gray!70, dashed, rounded corners=2pt, minimum width=81.9pt, minimum height=17.8pt, text=gray!60!black, align=center, font=\fontsize{5.50}{6.60}\selectfont\ttfamily] at (40.9pt,46.6pt) {\hyperref[thm:elemwiseProduct_has_vjp]{elemwiseProduct\_has\_vjp}\\(ch 1)};
  \node[draw=gray!70, dashed, rounded corners=2pt, minimum width=58.8pt, minimum height=17.8pt, text=gray!60!black, align=center, font=\fontsize{5.50}{6.60}\selectfont\ttfamily] at (121.9pt,46.6pt) {\hyperref[thm:identity_has_vjp]{identity\_has\_vjp}\\(ch 1)};
  \draw[->, >=stealth, dashed, gray!60, line width=0.40pt] (49.4pt,37.5pt) .. controls (54.5pt,32.7pt) and (61.0pt,26.4pt) .. (66.9pt,20.9pt) -- (75.1pt,13.0pt);
  \draw[->, >=stealth, dashed, gray!60, line width=0.40pt] (113.2pt,37.5pt) .. controls (108.1pt,32.7pt) and (101.5pt,26.6pt) .. (95.6pt,21.1pt) -- (87.2pt,13.2pt);
\end{tikzpicture}
\end{center}

% --- new:
%   \section{The theorems}
%   
%   \begin{center}
\begin{tikzpicture}[x=1pt,y=1pt, every node/.style={inner sep=1.50pt, font=\fontsize{7.00}{8.40}\selectfont\ttfamily}]
  \node[draw=green!50!black, fill=green!15, rectangle, rounded corners=4.0pt, align=center, minimum width=69.0pt, minimum height=12.7pt] at (80.9pt,6.4pt) {\hyperref[thm:seBlock_has_vjp]{seBlock\_has\_vjp}};
  \node[draw=gray!70, dashed, rounded corners=2pt, minimum width=81.9pt, minimum height=17.8pt, text=gray!60!black, align=center, font=\fontsize{5.50}{6.60}\selectfont\ttfamily] at (40.9pt,46.6pt) {\hyperref[thm:elemwiseProduct_has_vjp]{elemwiseProduct\_has\_vjp}\\(ch 1)};
  \node[draw=gray!70, dashed, rounded corners=2pt, minimum width=58.8pt, minimum height=17.8pt, text=gray!60!black, align=center, font=\fontsize{5.50}{6.60}\selectfont\ttfamily] at (121.9pt,46.6pt) {\hyperref[thm:identity_has_vjp]{identity\_has\_vjp}\\(ch 1)};
  \draw[->, >=stealth, dashed, gray!60, line width=0.40pt] (49.4pt,37.5pt) .. controls (54.5pt,32.7pt) and (61.0pt,26.4pt) .. (66.9pt,20.9pt) -- (75.1pt,13.0pt);
  \draw[->, >=stealth, dashed, gray!60, line width=0.40pt] (113.2pt,37.5pt) .. controls (108.1pt,32.7pt) and (101.5pt,26.6pt) .. (95.6pt,21.1pt) -- (87.2pt,13.2pt);
\end{tikzpicture}
\end{center}


% --- old:
%   definition every chapter uses; a statement that leans on many others shows only
%   what fits, and its \texttt{\textbackslash uses} line has the rest; every box is
%   a link.}
% --- new:
%   definition every chapter uses; a statement that leans on many others shows only
%   what fits, and its \texttt{\textbackslash uses} line has the rest; one that
%   stands only on what every chapter uses hangs from the figure's shared box, which
%   says where those statements live; every box is a link.}

